\section{Simplified waste-generation schemes, and open items}\label{sec:mkt:schemes}

Section~\ref{sec:sp:schemes} listed eight simplifications of the waste-generation scheme and classified
them by whether they change the allocation or only the reported mass flows. The useful question here is
a different one: whether each changes the \emph{instrument set}. The two questions have different
answers, and the difference is worth having in one place, because a switch that changes the allocation
need not change what the government has to do, and a switch that leaves the allocation alone can still
remove an instrument.

Throughout, references to the switches are to Section~\ref{sec:sp:schemes}, and the instrument names
are those of Table~\ref{tab:mkt:instruments}.

\begin{enumerate}[label=\textbf{(W\arabic*)},ref=W\arabic*,leftmargin=3.2em,itemsep=0.35em]

\item \textbf{No nature-attributable waste} (\ref{sw:sp:nature}), $\Omega^{N,S}=\Omega^{D,S}=0$. The
two displacement charges $\tau^{N,S}$ and $\tau^{D,S}$ vanish, under both closures. This is the single
largest simplification of the menu: with it, \emph{no instrument at all is levied on the mine}, and the
extraction path is corrected entirely through the deposit that the material bears once it reaches
production. The mine's problem becomes the textbook Hotelling problem with an unaltered cost function,
and the whole of the materials correction sits on $R$, on $W$ and on $G$. Under Version~B the menu
reduces to three instruments plus the residual $s^W$.

\item \textbf{Waste from output only} (\ref{sw:sp:omegaY}), $\omega^Y=1$ and the rest zero. Under
Version~B nothing changes: the composition block is already empty. Under Version~A the block collapses
to a single ad valorem tax on output, $\tau^Y=\Delta$, and the gate fee of
Proposition~\ref{prop:mkt:A:gate} has base $\Omega Y$ --- a fee paid entirely by the final-good
producer.

\item \textbf{Waste from output and consumption only} (\ref{sw:sp:omegaYC}). Under Version~B, again
nothing. Under Version~A this is the smallest scheme in which the household faces any environmental
instrument at all: the consumption tax $\tau^C=\Delta\omega^C$ appears here and in no smaller scheme.

\item \textbf{Uniform intensities} (\ref{sw:sp:uniform}), $\omega^j=1$ for every $j$. Under Version~B,
nothing. Under Version~A the five composition instruments become one uniform rate $\Delta$, levied with
a plus sign on every use of the final good and on output; equivalently, a gate fee $\zeta$ on total
waste attributable to final-good spending, which is the cleanest form the composition block ever takes
and the one a quantitative implementation should use first.

\item \textbf{Capital embodies nothing} (\ref{sw:sp:phiI}), $\Phi^I=0$ or $\phi^I=0$. The capital
instrument $\tau^K$ vanishes, $Q=1$, and the price $p^M$ disappears with the state $M^K$. Under
Version~A, $\chi=0$ kills $\zeta$ and $\Delta$, so the entire composition block vanishes as well and
the two closures require the \emph{identical} instrument set:
\begin{align}
\tau^R &= p^W,
&
s^W &= p^W,
&
\tau^E &= p^P,
\notag\\
\tau^{N,S} &= \Omega^{N,S}p^W,
&
\tau^{D,S} &= \Omega^{D,S}p^W,
&
s^R &= d^Wp^P .
\label{eq:mkt:schemes:minimal}
\end{align}
This is the minimal menu of the model --- a source-neutral materials deposit--refund, an emission tax
with a recovery credit, and two displacement charges --- and, combined with~\ref{sw:sp:nature}, it
reduces to a deposit--refund and an emission tax alone. It is the natural first target for the
numerical implementation, for the same reason it was in Section~\ref{sec:sp:schemes}.

\item \textbf{No exploration} (\ref{sw:sp:explore}), $D\equiv0$. The instruments $\tau^{D,S}$ and
$\tau^D$ go with it, and the mine holds one stock rather than two.

\item \textbf{Exploration costs independent of past discovery} (\ref{sw:sp:CDX}), $C^D_X=0$. The mine
stops holding $X$ and $P^X\equiv0$, but the instrument set is unchanged. What changes is the behaviour
of the instruments over time: with a genuine steady state restored, every instrument in
Table~\ref{tab:mkt:instruments} converges to a constant, computable from the quasi-stationary
block~\eqref{eq:sp:B:quasistationary}. This is the switch to use when the object of interest is the
long-run level of the optimal policy rather than its transition path.

\item \textbf{No pollution channel in production} (\ref{sw:sp:nopoll}), $F_P=0$. The instrument set is
unchanged in form and only the level of $p^P$ moves. One thing does change qualitatively: with $F_P=0$
the firms suffer no private loss from pollution, so the emission tax carries the entire correction,
whereas with $F_P<0$ part of the damage is already borne privately through lower output --- though not
internalised, since $P$ is taken as given.

\end{enumerate}

\begin{table}[H]
\centering
\caption{What each simplification does to the optimal instrument set.}
\label{tab:mkt:switches}
\setlength{\tabcolsep}{5pt}
\renewcommand{\arraystretch}{1.05}
\begin{tabular}{@{}p{0.30\textwidth}p{0.30\textwidth}p{0.32\textwidth}@{}}
\toprule
Switch & Version B & Version A \\
\midrule
\ref{sw:sp:nature}\ \ no $\mathcal N$        & removes $\tau^{N,S},\tau^{D,S}$ & removes $\tau^{N,S},\tau^{D,S}$ \\
\ref{sw:sp:omegaY}\ \ waste from $Y$ only    & no change (block empty) & block collapses to $\tau^Y$ \\
\ref{sw:sp:omegaYC}\ \ waste from $Y,C$ only & no change & block collapses to $\tau^Y,\tau^C$ \\
\ref{sw:sp:uniform}\ \ uniform $\omega^j$    & no change & block becomes one uniform rate $\Delta$ \\
\ref{sw:sp:phiI}\ \ $\Phi^I=0$               & removes $\tau^K$ & removes $\tau^K$ and the whole block; menus coincide \\
\ref{sw:sp:explore}\ \ no exploration        & removes $\tau^{D,S}$ & removes $\tau^{D,S},\tau^D$ \\
\ref{sw:sp:CDX}\ \ $C^D_X=0$                 & unchanged; instruments converge & unchanged; instruments converge \\
\ref{sw:sp:nopoll}\ \ $F_P=0$                & unchanged in form & unchanged in form \\
\bottomrule
\end{tabular}
\end{table}

Two readings of the table matter. First, the simplifications that are free under Version~B --- the
composition switches~\ref{sw:sp:omegaY}--\ref{sw:sp:uniform}, inert for the allocation by
Proposition~\ref{prop:sp:inert} --- are equally free for the policy, and for the same reason. Under
Version~A they are neither: they change both the allocation and the number of instruments needed to
support it. Second, no switch removes the deposit--refund pair. Mass conservation is not a modelling
option, and no setting of the parameters produces a version of the model in which the charge on
materials can be dispensed with or replaced by a charge on disposal.

\begin{remark}[The materials block does not balance]\label{rem:mkt:revenue}
The deposit--refund scheme is not revenue-neutral, and the imbalance is informative. Under Version~B
with $\eta=0$, the net revenue of the materials block is
\begin{align}
\underbrace{p^W\left(R+\mathcal N\right)}_{\text{deposits}}
-\underbrace{p^WW}_{\text{refunds}}
-\underbrace{\Phi^I\left(p^W-p^M\right)G}_{\text{capital rebate}}
&= \Phi^Ip^MG-p^W\delta M^K,
\label{eq:mkt:schemes:revenue}
\end{align}
using the ledger $W=R-\dot M^K+\mathcal N$ and $\dot M^K=\Phi^IG-\delta M^K$. The government accrues
$p^M$ per tonne of mass placed in the durable stock and pays out $p^W$ per tonne released from it: the
materials block is a liability account for deposits not yet refunded. In a materially stationary
economy, $\Phi^IG=\delta M^K$, and~\eqref{eq:mkt:schemes:revenue} is $\delta M^K\left(p^M-p^W\right)<0$
whenever $p^W>p^M$ --- the block runs a permanent deficit, financed lump-sum, equal to the interest on
the deposits embodied in the capital stock. Emission tax revenue $p^P\Xi\ge0$ is separate and does not
offset it systematically.
\end{remark}


\subsection{Open items specific to the market economy}\label{sec:mkt:open}

These are additional to the six items of Section~\ref{sec:sp:ext:open}, all of which carry over: in
particular the implementation propositions above invoke uniqueness of the planner's
optimum~(\ref{op:sp:soc}), and the corner behaviour of~\ref{op:sp:corner} is inherited.

\begin{enumerate}[label=\textbf{M\arabic*.},ref=M\arabic*,leftmargin=2.6em,itemsep=0.35em]

\item\label{op:mkt:dumping} \textbf{Enforcement of the collection obligation.} The waste-management
sector is assumed to collect all waste, and generators are assumed to hand it over. Neither is free.
If generators can evade collection, the equilibrium changes qualitatively, and in a direction that
strengthens the case for the entry-side charge: a deposit levied on material entering production
cannot be evaded by disposing of the waste elsewhere, whereas a gate fee can, and a gate fee is
precisely the instrument that creates the incentive to evade. This is a standard argument for
deposit--refund schemes and it is orthogonal to the one made in
Sections~\ref{sec:mkt:B:nogate} and~\ref{sec:mkt:A:gate}, which holds with perfect enforcement. It is
not modelled here.

\item\label{op:mkt:corner} \textbf{The mass constraint has no market
(Remark~\ref{rem:mkt:B:eta}).} Under Version~B the restriction $\Phi^IG\le R$ binds on no private
agent. The multiplier $\eta$ has a market representation --- a uniform reduction of the deposit and an
offset to the capital rebate --- but whether a decentralised equilibrium exists at the corner, and what
happens to it under a policy that does not include $\eta$, is not established. Under Version~A the
question does not arise.

\item\label{op:mkt:secondbest} \textbf{Welfare under an arbitrary policy.} The equilibrium system of
Sections~\ref{sec:mkt:A:system} and~\ref{sec:mkt:B:information} determines the allocation for any
instrument path, which is what a status-quo or partial-policy counterfactual needs and which
Part~\ref{part:sp} does not provide. What is still missing is the evaluation of $U_0$ along such a
path and therefore the welfare ranking of two policies; this is~\ref{op:sp:nonopt}, and this part
resolves only its positive half. The natural second-best exercises --- what to do when $\tau^K$ is
unavailable, or when the deposit must be levied on virgin material only --- follow immediately from
the same system and have not been carried out.

\item\label{op:mkt:existence} \textbf{Existence and uniqueness of equilibrium.} Nothing here
establishes that the equilibrium of Definition~\ref{def:mkt:equilibrium} exists or is unique for given
instrument paths, and under Version~A the accounting block is solved by the positive root of a
quadratic whose coefficients depend on the allocation, so the fixed point is not obviously well
behaved. The transfer $T$ is also assumed lump-sum and unrestricted in sign, which
Remark~\ref{rem:mkt:revenue} shows is not innocuous: the materials block runs a persistent deficit in a
materially stationary economy.

\item\label{op:mkt:tender} \textbf{The collection tender.} The result $s^W=p^W$ rests on the
collection sector being competitive and earning exactly zero profit. With market power, or with a
tender that leaves rents, the break-even payment and the shadow price of waste part company, and the
identification of the two --- the cleanest result in this part --- is lost. Since waste collection is
close to a natural monopoly in practice, this is a substantive limitation rather than a technical one.

\item\label{op:mkt:instrumentsign} \textbf{Instruments of ambiguous sign.} By~\eqref{eq:sp:B:pW-sign}
and~\eqref{eq:sp:A:pW-sign}, $p^W$ turns negative in the late transition, at which point the deposit
becomes a payment for drawing material in, the refund a charge for handing waste over, and the capital
rebate a tax. Whether such a scheme is implementable in the sense the model assumes --- in particular
whether a negative refund is enforceable, which is the same enforcement question as~\ref{op:mkt:dumping}
with the sign reversed --- is not addressed.

\end{enumerate}
